% Guion de entrevista semiestructurada - Encargado del Gimnasio Universitario (UBU)
\begin{table}[H]
  \centering\small
  \caption{Datos generales de la entrevista.}\label{tab:datos-entrevista}
  \begin{tabularx}{\textwidth}{@{}lL@{}}
    \toprule
    Campo & Detalle \\
    \midrule
    Proyecto & \sistema \\
    Entrevistado & Encargado del Gimnasio Universitario (UBU) -- PI-02 \\
    Tipo & Semiestructurada \\
    Objetivo & Conocer el proceso actual de inscripción, asignación de horarios, control de asistencia, registro de información de salud y manejo de equipos del gimnasio, para identificar necesidades del sistema. \\
    Duración estimada & 45 minutos \\
    Entrevistadores & \pendiente{integrantes y quién toma notas} \\
    Registro & Notas de campo y grabación de audio solo con consentimiento (Anexo~\ref{anx:consentimiento}) \\
    \bottomrule
  \end{tabularx}
\end{table}

\subsection*{Apertura}
\begin{enumerate}[label=\alph*)]
  \item Presentación del equipo y del propósito de la entrevista.
  \item Lectura y firma del consentimiento informado; solicitud de permiso para grabar.
  \item Indicación de la duración aproximada y de que puede omitir cualquier pregunta.
\end{enumerate}

\subsection*{Contexto y rol}
\begin{enumerate}[label=\arabic*., series=preguntas]
  \item ¿Qué funciones cumple usted dentro del Gimnasio Universitario?
  \begin{itemize}
    \item[R:] ¿Desde hace cuánto tiempo está a cargo de esta área?
    \item[R:] ¿Quién más trabaja en los distintos turnos?
  \end{itemize}
\end{enumerate}

\subsection*{Bloque 1: Proceso de inscripción (trámite en TICs y ficha UBU)}
\begin{enumerate}[resume*=preguntas]
  \item Cuéntenos paso a paso, ¿cómo hace hoy un estudiante desde que quiere entrar al gimnasio hasta que ya puede entrenar?
  \begin{itemize}
    \item[R:] ¿Por qué tienen que ir primero a TICs por la hoja?
    \item[R:] ¿Qué pasa después de que se la entregan a usted?
    \item[R:] ¿Pueden inscribirse solo estudiantes o también docentes y administrativos?
  \end{itemize}
  \item ¿Cómo describe su experiencia actual gestionando las inscripciones mediante las fichas impresas?
  \begin{itemize}
    \item[R:] ¿Qué inconvenientes se han presentado con las fichas?
  \end{itemize}
  \item En la ficha hay un espacio de «Observaciones» para uso institucional, ¿qué datos suelen anotar ahí?
  \begin{itemize}
    \item[R:] ¿Qué tipo de informes o estadísticas le pide Bienestar Universitario y cada cuánto tiempo?
  \end{itemize}
\end{enumerate}

\subsection*{Bloque 2: Horarios, cupos y control de asistencia}
\begin{enumerate}[resume*=preguntas]
  \item ¿Cómo se determina y distribuye la capacidad de usuarios en cada uno de los siete bloques horarios diarios?
  \begin{itemize}
    \item[R:] ¿El último turno de la tarde (17h10 a 17h50) se maneja igual que los demás aunque dure menos tiempo?
  \end{itemize}
  \item ¿Qué ocurre cuando varios estudiantes marcan el mismo horario y se supera el límite de espacio?
  \begin{itemize}
    \item[R:] ¿Cómo se entera el estudiante si ya no hay cupo en la hora que eligió?
  \end{itemize}
  \item ¿De qué manera controlan hoy quién asiste cada día a entrenar?
  \begin{itemize}
    \item[R:] ¿Qué pasa si alguien se inscribe pero deja de venir?
    \item[R:] ¿Existe un número de faltas a partir del cual se pierde el cupo?
  \end{itemize}
  \item ¿Cómo verifican en la puerta que la persona que entra sí esté inscrita en ese turno y no en otro?
  \begin{itemize}
    \item[R:] ¿Qué hacen si llega alguien que no está inscrito o quiere entrenar fuera de su horario?
  \end{itemize}
\end{enumerate}

\subsection*{Bloque 3: Salud, lesiones y rutinas}
\begin{enumerate}[resume*=preguntas]
  \item ¿Se realiza alguna evaluación física o se pide información previa cuando un estudiante nuevo empieza a entrenar?
  \begin{itemize}
    \item[R:] ¿Llevan algún registro del peso, la estatura o los avances de cada persona?
  \end{itemize}
  \item Cuando un estudiante indica en la ficha que tiene una lesión o condición médica, ¿qué se hace con esa información?
  \begin{itemize}
    \item[R:] ¿Se le restringe algún ejercicio o equipo?
    \item[R:] ¿Qué datos le gustaría tener a la mano si ocurre una emergencia durante el entrenamiento?
  \end{itemize}
  \item ¿De qué manera consiguen los estudiantes la rutina de ejercicios que van a realizar?
  \begin{itemize}
    \item[R:] Cuando un estudiante no sabe usar una máquina, ¿a quién acude?
  \end{itemize}
  \item ¿Cómo manejan el estado de las máquinas y equipos cuando alguno se daña o necesita mantenimiento?
  \begin{itemize}
    \item[R:] ¿Cómo avisan a los usuarios o a Bienestar Universitario sobre los equipos fuera de servicio?
  \end{itemize}
\end{enumerate}

\subsection*{Cierre: prioridades y viabilidad}
\begin{enumerate}[resume*=preguntas]
  \item De todas sus actividades diarias en el gimnasio, ¿cuál considera que le toma más tiempo o esfuerzo?
  \begin{itemize}
    \item[R:] ¿Qué le gustaría que este sistema resuelva primero?
  \end{itemize}
  \item ¿Qué autoridades o dependencias deben aprobar el uso de un sistema como este?
  \item ¿Hay algún tema que no hayamos tocado y que considere importante?
  \item ¿Con qué otra persona nos recomienda conversar sobre el funcionamiento del gimnasio?
\end{enumerate}

\begin{enumerate}[label=\alph*), start=4]
  \item Resumen de lo entendido y validación con el entrevistado.
  \item Agradecimiento y acuerdo sobre el siguiente contacto.
\end{enumerate}
